\documentclass[UTF8, a4paper, zihao=-4]{ctexart}
\usepackage{geometry}
\geometry{margin=2.5cm}
\usepackage{booktabs}
\usepackage{graphicx}
\usepackage{url}
\usepackage{hyperref}
\usepackage{listings}
\lstset{basicstyle=\ttfamily\small, breaklines=true, frame=single}

\title{\textbf{文献综述智能体设计说明}\\{\large 大模型智能体长期记忆综述项目}}
\author{成员 A\quad 成员 B\quad 成员 C}  % 待真人姓名替换
\date{2026 年 10 月}

\begin{document}
\maketitle

\section{设计目标与约束}

本智能体的任务：围绕研究问题``大模型智能体如何从持续交互中保留有用经验，
在信息变化时正确更新记忆，并在后续任务中可靠地使用这些记忆''，完成一篇
不少于 30 篇真实文献的综述的全部可自动化环节——检索、去重、筛选、证据
提取、陈述登记与草稿生成。

核心约束（对应课程要求）：
\begin{itemize}
  \item \textbf{真实性优先}：任何引用、摘录、数值必须可回查到原文；
        禁止编造。宁可输出``待核验''，不输出无依据结论。
  \item \textbf{可恢复}：跨运行保留研究状态，中断后可继续且不重复写入。
  \item \textbf{人工节点}：真人核验不可替代；智能体只准备核验材料，
        不代签、不预填``已通过''。
\end{itemize}

\section{架构}

\begin{lstlisting}
研究目标 / 当前状态 / 人工反馈
              |
协调 Agent（coordinator.py）
  观察状态(RQ 覆盖/证据缺口) -> 选择工具 -> 执行 -> 记录
              |
工具层（确定性为主，模型负责规划/提取/综合）
  search: OpenAlex / arXiv / Crossref / Semantic Scholar
  dedup:  DOI / arXiv ID / 题名-作者-年份 三重键归并 -> work_id
  fetch:  PDF 下载(sha256 存档) + pdftotext 页级提取 + HTML 通道
  evidence/claims/memory/review: JSONL 事件溯源读写
  checks: 引用解析 / 证据支持 / 语料计数
              |
持久化（data/*.jsonl + runs/<id>/events.jsonl，全部追加式）
              |
带引用草稿 + 缺口清单 + 人工核验包
\end{lstlisting}

两层运行形态（决策记录 \texttt{docs/decisions.md} D-2）：
\begin{itemize}
  \item \textbf{协调器回合}（G2 演示）：模型在工具循环中自主决策——
        观察缺口、选择检索词、读页、提交证据、诚实报告失败；
  \item \textbf{批处理管线}（G3 规模化）：逐篇单次模型调用提取证据，
        与协调器共用同一套摘录验证器与数据格式。
\end{itemize}

\section{记忆设计}

四类记忆（\texttt{memory\_events.jsonl}，追加事件式）：
任务记忆（范围/协议/待办）、证据记忆（定位/摘录/解释）、
决策记忆（筛选理由/修改原因）、经验记忆（流程总结，默认未经事实核验）。

两类状态分离记录：
\begin{itemize}
  \item \texttt{valid\_status}: active / superseded / disputed / rejected——
        有效性；默认检索过滤非 active；
  \item \texttt{trust\_status}: unverified / source\_checked / human\_verified——
        核验层级；草稿与模型推测不得自动升级。
\end{itemize}

更新采用 \texttt{supersedes} 事件链：新条目取代旧条目但历史完整保留，
可审计``旧错误经摘要后失去来源''的路径。冲突不按时间新旧自动覆盖：
研究范围变更可更新任务规则；两篇论文结论不同则保留各自实验条件并标记差异。

\section{反幻觉机制与实测}

\subsection{机制}

摘录提交时的\textbf{逐字签名验证}（代码强制，不依赖模型自觉）：
\begin{enumerate}
  \item 摘录与页面提取文本各生成``签名''——剔除空白与连字符（排版产物）
        后的字符序列；
  \item 摘录签名必须是页面签名的\textbf{连续子串}；
  \item 通过后回存\textbf{原文真实片段}（保留原始连字符与换行），
        保证 \texttt{evidence.excerpt} 与提取文本逐字一致；
  \item 失败则拒绝并返回页面上最相似的原文片段（相似度提示），
        供模型自纠错重试。
\end{enumerate}

改写、缩写、省略号压缩、凭记忆复述、表格单元格拼接均会被拒绝。

\subsection{实测统计（真实运行日志汇总）}

\begin{table}[h]
\centering
\begin{tabular}{lrr}
\toprule
运行形态 & 摘录被拒 & 摘录接受 \\
\midrule
协调器回合（G2，12 次运行） & 258 & 38 \\
批处理管线（G3 深读） & 102 & 313 \\
\bottomrule
\end{tabular}
\caption{逐字验证拦截统计（真实运行日志汇总，\texttt{runs/}）。协调器
早期回合的高拒绝率来自验证规则迭代（连字符/双栏问题修复前）与早期提示词
不成熟；被拒摘录从未进入证据库。独立抽查：随机 4 条证据回查原 PDF 页，
4/4 摘录逐字存在。}
\end{table}

协调器回合中两次典型的真实决策（\texttt{runs/} 日志可回查）：
\begin{itemize}
  \item \textbf{缺口驱动}：模型观察到 RQ1/RQ5 证据不足后，自主选择 7 个
        补充检索词（如 ``long-term memory versus context window RAG
        parametric knowledge boundary''），非脚本预设；
  \item \textbf{诚实失败}：工具链缺少证据枚举能力时，模型拒绝编造
        evidence\_id，明确报告``在禁止编造约束下无法合规完成''并提出
        工具改进建议（后被采纳实现）。
\end{itemize}

\section{人工核验节点}

\begin{itemize}
  \item 核验包生成器（\texttt{review\_queue}）从真实数据产生三层待办：
        文献身份/版本核验（每项候选引用文献）、边界判定复核（模型纳入/
        排除决策）、claims 复核（direct/synthesis 全部）；
  \item 核验结果由真人写入 \texttt{data/reviews.jsonl}，
        智能体不预填``已通过''记录；
  \item 优先级标注：更新与可靠性分支（综述重点）为 P0。
\end{itemize}

\section{评估设计}

12 个冻结案例（\texttt{evaluation/cases.jsonl}）× 三条件
（A 无记忆基线 / B 上下文注入 / C 结构化记忆）：
跨会话恢复、失效记忆处理、不确定性保留、跨运行任务四类各 3 例。
评分规则见 \texttt{evaluation/rubric.md}。

\subsection{结果（v6：三态判分——成功/失败/待核验）}

\begin{table}[h]
\centering
\begin{tabular}{lccc}
\toprule
 & A（无记忆） & B（上下文注入） & C（结构化记忆） \\
\midrule
成功 & 2/12 & \textbf{11}/12 & 9/12 \\
失败 & 10/12 & 0/12 & 3/12 \\
待核验 & 0 & 1 & 0 \\
\bottomrule
\end{tabular}
\caption{任务完成度三态统计。multirun 为真实两阶段实验；M4-03 开放式
内容主张按保守规则转人工核验（D-13），不默认成功。}
\end{table}

\textbf{结论}：B（上下文注入）成功率最高（11 成功/1 待核验）；C（结构化
记忆）9 成功 3 失败（动作缺失与过度确定），跨运行持久化经两阶段实验真实验证；
A（无记忆）2 成功且以编造补偿。开放式内容主张的相符性判断超出自动规则
边界，已转人工核验——自动检查不冒充语义验证。

\subsection{过程完整性}

评估经历两轮实现修复（记忆检索对中文失效、失效记忆过滤缺陷），
每次修复后全量重跑并在 \texttt{docs/decisions.md} D-4/D-5 记录；
受污染的早期结果归档为 \texttt{results\_*.bak}，未计入报告。

\section{局限}

\begin{itemize}
  \item 摘录验证保证``摘录逐字来自所指页面''，不保证``解释无过度引申''——
        语义层核验仍需真人；
  \item 页级定位基于 pdftotext raw 模式的页序，与印刷页码可能不一致
        （预印本无印刷页码）；核验包中保留原文 PDF 供人工对页；
  \item 批处理管线的分支标注与矩阵分类由模型判定，属待复核信息
        （已列入核验队列）；
  \item 评估为小样本演示性评估（12 案例 × 3 条件），不构成方法优劣的
        统计证据。
\end{itemize}

\end{document}
